\documentclass[12pt,a4paper]{article}

\usepackage{fontspec}
\usepackage{polyglossia}
\setdefaultlanguage{russian}
\setotherlanguage{english}
\setmainfont{Arial}
\setsansfont{Arial}
\setmonofont{Arial}
\newfontfamily\cyrillicfonttt{Arial}
\usepackage[left=2cm,right=2cm,top=2cm,bottom=2cm]{geometry}
\usepackage{graphicx}
\usepackage{amsmath}
\usepackage{booktabs}
\usepackage{array}
\usepackage{float}
\usepackage{xcolor}
\usepackage{hyperref}
\usepackage{xurl}
\usepackage{fancyhdr}
\usepackage{microtype}

\definecolor{accent}{HTML}{2864A8}
\definecolor{textcolor}{HTML}{243B53}
\hypersetup{colorlinks=true,linkcolor=accent,urlcolor=accent}
\color{textcolor}
\pagestyle{fancy}
\fancyhf{}
\lhead{Поиск частых наборов --- Apriori}
\rhead{\thepage}
\setlength{\headheight}{15pt}
\setlength{\parindent}{1.25cm}
\setlength{\parskip}{0.35em}

\begin{document}

\begin{titlepage}
  \centering
  \vspace*{3cm}
  {\Huge\bfseries ОТЧЁТ О ВЫПОЛНЕНИИ ЗАДАНИЯ\par}
  \vspace{1cm}
  {\LARGE\color{accent}\bfseries Поиск частых наборов объектов\\алгоритмом Apriori\par}
  \vspace{3cm}
  \fbox{\begin{minipage}{0.78\textwidth}
    Набор данных: \textbf{baskets.csv}\\
    Число транзакций: \textbf{7501}\\
    Исследованные пороги: \textbf{1\%, 3\%, 5\%, 10\%, 15\%}
  \end{minipage}}
  \vfill
  {2026\par}
\end{titlepage}

\section{Формулировка задания}

Требуется разработать программу поиска частых наборов объектов в транзакционном наборе данных методом Apriori или его модификацией. Для каждого найденного набора необходимо указать поддержку. Входными параметрами служат набор данных, минимальный порог поддержки и способ сортировки результата: по убыванию поддержки или лексикографически.

Экспериментальная часть должна сравнивать время работы при разных порогах поддержки и показывать число частых наборов разной длины. Результаты необходимо представить диаграммами и объяснить наблюдаемые зависимости.

\section{Материалы проекта}

Исходный код, набор данных и результаты размещены в репозитории:
\url{https://github.com/FeodorG/data-mining-fundamentals}.

\section{Методика}

Реализован классический поуровневый Apriori. На первом шаге подсчитывается поддержка отдельных товаров. Для каждого следующего уровня выполняются объединение частых наборов предыдущей длины, отсечение кандидатов с нечастыми подмножествами и один проход по транзакциям для подсчёта поддержки кандидатов.

Относительная поддержка набора $X$ вычисляется как
\[
  \operatorname{support}(X)=\frac{|\{t:X\subseteq t\}|}{N},
\]
где $N$ --- число транзакций. Абсолютный порог равен округлённому вверх произведению относительного порога на $N$. Время каждого опыта измерялось трижды; далее приведена медиана.

Программа автоматически читает UTF-8 и Windows-1251. Результат сохраняется в CSV с набором, его длиной, абсолютной и относительной поддержкой. Параметр \texttt{--order} принимает значения \texttt{support} и \texttt{lexicographic}.

\section{Результаты экспериментов}

\begin{table}[H]
\centering
\caption{Сводные результаты, три запуска для каждого порога}
\begin{tabular}{rrrr}
\toprule
Порог & Медиана, с & Частых наборов & Макс. длина \\
\midrule
1\%  & 1.2628 & 261 & 3 \\
3\%  & 0.2019 & 55  & 2 \\
5\%  & 0.1200 & 28  & 2 \\
10\% & 0.0130 & 7   & 1 \\
15\% & 0.0087 & 5   & 1 \\
\bottomrule
\end{tabular}
\end{table}

\begin{figure}[H]
  \centering
  \includegraphics[width=0.95\textwidth]{results/runtime.png}
  \caption{Зависимость медианного времени работы от порога поддержки}
\end{figure}

\section{Распределение наборов по длине}

\begin{figure}[H]
  \centering
  \includegraphics[width=0.95\textwidth]{results/itemsets_by_length.png}
  \caption{Количество частых наборов; цвет показывает длину набора}
\end{figure}

\begin{table}[H]
\centering
\caption{Число наборов каждой длины}
\begin{tabular}{rrrrr}
\toprule
Порог & Длина 1 & Длина 2 & Длина 3 & Всего \\
\midrule
1\%  & 74 & 170 & 17 & 261 \\
3\%  & 36 & 19  & 0  & 55  \\
5\%  & 25 & 3   & 0  & 28  \\
10\% & 7  & 0   & 0  & 7   \\
15\% & 5  & 0   & 0  & 5   \\
\bottomrule
\end{tabular}
\end{table}

\section{Интерпретация}

При повышении порога множество допустимых наборов монотонно сокращается: с 261 при 1\% до 5 при 15\%. Это прямое следствие определения частого набора: более строгому порогу удовлетворяет меньше товаров и комбинаций.

На пороге 1\% преобладают пары (170 из 261), а также появляются 17 троек. При 3\% остаются 19 пар, при 5\% --- только 3 пары. Начиная с 10\% частыми являются лишь одиночные товары. Следовательно, совместные покупки имеют заметно меньшую поддержку, чем самые популярные отдельные товары.

Низкий порог сохраняет больше наборов на каждом уровне, поэтому Apriori генерирует и проверяет существенно больше кандидатов. При высоких порогах процесс завершается уже после уровня одиночных товаров. Небольшие колебания времени между запусками объясняются состоянием операционной системы.

\section{Вывод}

Разработанная программа выполняет поиск частых наборов, рассчитывает абсолютную и относительную поддержку и поддерживает оба требуемых порядка сортировки. Эксперименты показывают ожидаемый компромисс: низкий порог позволяет выявить больше сложных сочетаний, но увеличивает вычислительные затраты; высокий порог оставляет только наиболее массовые товары и резко ускоряет поиск.

\end{document}
